% Appendix: Isolation Methods — Study Levers
% Standalone mathematical reference for the design space of isolation methods.
\PassOptionsToPackage{unicode}{hyperref}
\PassOptionsToPackage{hyphens}{url}
\documentclass[
]{article}
\usepackage{xcolor}
\usepackage{amsmath,amssymb}
\setcounter{secnumdepth}{-\maxdimen} % remove section numbering
\usepackage{iftex}
\ifPDFTeX
  \usepackage[T1]{fontenc}
  \usepackage[utf8]{inputenc}
  \usepackage{textcomp} % provide euro and other symbols
\else % if luatex or xetex
  \usepackage{unicode-math} % this also loads fontspec
  \defaultfontfeatures{Scale=MatchLowercase}
  \defaultfontfeatures[\rmfamily]{Ligatures=TeX,Scale=1}
\fi
\usepackage{lmodern}
% Use upquote if available, for straight quotes in verbatim environments
\IfFileExists{upquote.sty}{\usepackage{upquote}}{}
\IfFileExists{microtype.sty}{% use microtype if available
  \usepackage[]{microtype}
  \UseMicrotypeSet[protrusion]{basicmath} % disable protrusion for tt fonts
}{}
\makeatletter
\@ifundefined{KOMAClassName}{% if non-KOMA class
  \IfFileExists{parskip.sty}{%
    \usepackage{parskip}
  }{% else
    \setlength{\parindent}{0pt}
    \setlength{\parskip}{6pt plus 2pt minus 1pt}}
}{% if KOMA class
  \KOMAoptions{parskip=half}}
\makeatother
\setlength{\emergencystretch}{3em} % prevent overfull lines
\providecommand{\tightlist}{%
  \setlength{\itemsep}{0pt}\setlength{\parskip}{0pt}}
\usepackage{bookmark}
\IfFileExists{xurl.sty}{\usepackage{xurl}}{} % add URL line breaks if available
\urlstyle{same}
\hypersetup{
  hidelinks}

\newcommand{\dmodel}{d_{\mathrm{model}}}
\newcommand{\dhead}{d_{\mathrm{head}}}

\author{}
\date{}

\begin{document}

\section{Appendix: Isolation Methods — Study Levers}\label{isolation-methods-study-levers}

An \emph{isolation method} is a procedure that isolates the parts of the model or
residual stream responsible for learning a simple function in context, as a
lower-dimensional representation (a vector, a set of heads, a subspace) that can be
injected back into the model to induce the task. This appendix defines the design
space of such methods as four largely independent levers: the isolation algorithm,
the metric it optimises, the data it is fit on, and the metric by which it is judged.
The model is GPT-J-6B unless stated otherwise.

Motivating problem: zero-shot steering on held-out tasks is bimodal. Several tasks
(e.g.~synonym, word-length, country--currency) get essentially no steering benefit
from any current isolation product, while others steer at their few-shot ceiling. The
levers below define the design space; the hypotheses at the end define what a result
would mean.

\subsection{Notation}\label{notation}

\begin{itemize}
\tightlist
\item
  \(\dmodel\) is the residual-stream width and \(\dhead\) is the per-head width
  (\(\dmodel = 4096\), \(\dhead = 256\) for GPT-J, which has 28 layers of 16 heads,
  448 heads in total).
\item
  \(\mathcal{T}\) is the task universe and \(A \in \mathcal{T}\) is a task.
\item
  \(p_A^j\) is the \(j\)-th prompt for task \(A\), and
  \(\mathcal{P}_A = \{p_A^j\}_j\) is the prompt set for \(A\) (varying in ICL context
  length from 1 to 10 demonstrations, unless explicitly mentioned otherwise).
\item
  The ``:'' token is the final prompt token, i.e.~the pre-label position at which the
  model must produce the answer. All activations are read, and all isolation products
  are injected, at this position.
\item
  \(h(p_A^j) \in \mathbb{R}^{\dmodel}\) is the output of head \(h\) at the ``:''
  token on prompt \(p_A^j\), mapped into the residual stream.
\item
  The head vector \(h_A = \frac{1}{|\mathcal{P}_A|} \sum_j h(p_A^j)\) is the average
  head activation for task \(A\), taken over prompts the model answers correctly
  in context.
\item
  \(v_A = \sum_{h \in H} h_A\) is the function vector (FV) for task \(A\), where
  \(H\) is a selected subset of heads.
\item
  \(v^j_A = \sum_{h \in H} h(p_A^j)\) is the per-prompt function vector for prompt
  \(j\) on task \(A\).
\item
  \(z^t_\ell\) is the residual stream at layer \(\ell\) at token \(t\).
\end{itemize}

Every isolation product below is a vector (or family of vectors) injected additively
into the residual stream at the ``:'' token:
\[
z^{\text{``:''}}_\ell \;\leftarrow\; z^{\text{``:''}}_\ell + v. \tag{1}
\]

\subsection{Lever 1 — Isolation algorithm}\label{lever-1-isolation-algorithm}

Three algorithms are in play. Each produces, for a task \(A\), an injectable product.

\subsubsection{1a. CIE top-\texorpdfstring{\(k\)}{k} hard head selection}\label{lever-1a}

The Function Vectors algorithm of Todd et al. For every attention head \(h\), compute
the \emph{causal indirect effect} (CIE): patch the head's task-mean vector \(h_A\)
into prompts drawn from the chosen optimisation setting (Lever 2) and measure the
resulting improvement in recovery of the correct answer. Writing \(\tilde p\) for a
prompt from that setting and \(y\) for the correct answer,
\[
\mathrm{CIE}_A(h) \;=\; \mathbb{E}_{\tilde p}\!\left[\,
  p\big(y \mid \tilde p;\ h \leftarrow h_A\big) \;-\; p\big(y \mid \tilde p\big)
\,\right], \tag{2}
\]
where \(h \leftarrow h_A\) denotes replacing the head's output at the ``:'' token by
its task mean. The per-head score is averaged over the fitting tasks,
\(\mathrm{CIE}(h) = \frac{1}{|\mathcal{T}_{\mathrm{fit}}|}
\sum_{A \in \mathcal{T}_{\mathrm{fit}}} \mathrm{CIE}_A(h)\), and the shared head set
\(H\) is the hard selection of the top-\(k\) heads (canonically \(k = 40\)). The
isolation product is the function vector
\[
v_A \;=\; \sum_{h \in H} h_A, \tag{3}
\]
one vector per task, injected at a single layer (either swept over all layers with
the best layer reported, or a fixed convention layer). In the original formulation
the optimisation setting is the same-task shuffled-label 10-shot prompt; that is one
point of Lever 2, not part of the algorithm.

The selection criterion is \emph{implicit}: per-head answer recovery in a corrupted
(shuffled-label) 10-shot context. In its original form the algorithm never sees a
zero-shot prompt.

\subsubsection{1b. Sparse optimisation over heads}\label{lever-1b}

The algorithm of Hu et al.~2025 (arXiv:2505.05145, \S3.1), stated generically. Learn
a coefficient vector \(c \in [0,1]^{448}\) over all heads. The candidate vector
\[
v_A(c) \;=\; \sum_{h} c_h\, h_A \tag{4}
\]
is injected once at the ``:'' token of prompts drawn from the chosen optimisation
setting (Lever 2), and \(c\) is fit by minimising the teacher-forced negative
log-likelihood of the correct answer plus an \(\ell_1\) sparsity penalty,
\[
\mathcal{L}(c) \;=\;
\mathbb{E}\!\left[-\log p\big(y \mid \tilde p;\ v_A(c)\big)\right]
\;+\; \lambda \,\lVert c \rVert_1, \tag{5}
\]
with \(\lambda\) chosen by leave-one-task-out cross-validation over the fitting
tasks, and \(y\) scored as the full label (all of its tokens), not only its first
token. The isolation product used downstream is the hard set
\(H = \{h : c_h > 0.8\}\), summed \emph{unweighted}; that is, exactly the
construction of Eq.~(3) with the head list swapped.

The selection criterion is \emph{explicit and differentiable}, and generic over
Lever 2. The instantiation used in our experiments fits under the zero-shot setting
and yields a 23-head set.

\subsubsection{1c. Per-layer mean ``:''-token activation}\label{lever-1c}

No selection at all. For each layer \(\ell\), the isolation product is the mean
residual stream at the ``:'' token over 10-shot prompts of task \(A\),
\[
m^\ell_A \;=\; \frac{1}{|\mathcal{P}_A|} \sum_j z^{\text{``:''}}_\ell\!\big(p^j_A\big), \tag{6}
\]
one vector per layer, each injected at its own layer \(\ell\). This is the ``average
hidden state'' baseline of Todd et al.: it carries everything the residual stream
carries at that point (task identity, format, position, token statistics), not just
head-mediated task content.

\subsubsection{Compatibility note}\label{compatibility-note}

Algorithms 1a and 1b are both generic over the fitting objective (Lever 2): the CIE
can score heads by answer recovery in any prompt setting, and the sparse objective
can be evaluated under any prompt setting. Existing runs pin particular choices: the
CIE top-40 selection was fit under the same-task shuffled-label 10-shot setting, and
the sparse 23-head selection under the zero-shot setting. Algorithm 1c has no
optimisation step and is a fixed point of the crossing (though the prompts its mean
is taken over could in principle also vary).

\subsection{Lever 2 — Metric to optimise (fitting objective)}\label{lever-2-metric-to-optimise}

The prompt setting the isolation algorithm is fit against (applies to both 1a and
1b; see the compatibility note):

\begin{enumerate}
\tightlist
\item
  \textbf{Zero-shot accuracy.} Inject at the ``:'' token of a bare
  \texttt{Q:\ x\textbackslash nA:} prompt with no demonstrations; the objective in
  practice is the teacher-forced negative log-likelihood of the full label, a
  differentiable surrogate for accuracy.
\item
  \textbf{Uplift on same-task shuffled-label 10-shot prompts.} Ten demonstrations
  from task \(A\) with labels shuffled within the prompt (the corrupted context of
  Todd et al.); uplift is steered accuracy minus unsteered accuracy on the same
  prompts.
\item
  \textbf{Uplift on mixed-task mixed-label 10-shot prompts.} Demonstrations drawn
  from multiple tasks with labels mixed. The precise construction (which tasks mix,
  whether labels are additionally shuffled, how uplift is defined) is left open and
  must be fixed before first use.
\end{enumerate}

\subsection{Lever 3 — Data to fit on}\label{lever-3-data-to-fit-on}

\begin{enumerate}
\tightlist
\item
  \textbf{Train/test task split.} Fit on 20 train tasks and evaluate on 9 held-out
  test tasks (landmark--country, word-length, capitalize-first-letter, synonym,
  lowercase-first-letter, product--company, capitalize, country--currency, antonym).
  This is the transfer setting: the isolation product is shared across tasks it has
  never seen.
\item
  \textbf{Task-specific.} Fit on the evaluated task's own data. No generalisation
  claim; upper-bounds what the algorithm can do when allowed to see the task.
\item
  \textbf{All tasks pooled.} Fit on all 29 tasks together with no held-out tasks;
  measures the capacity of a \emph{shared} isolation product, not transfer.
\end{enumerate}

\subsection{Lever 4 — Metric of success (evaluation)}\label{lever-4-metric-of-success}

The same three metrics as Lever 2, now used as evaluations, crossed with \emph{where}
they are measured: on the train tasks, on the held-out tasks, or on the fitting task
itself. Two readout conventions coexist and are \textbf{not numerically comparable}:

\begin{itemize}
\tightlist
\item
  \textbf{Best-layer first-token readout.} First-token top-1 accuracy of the gold
  answer; the injection layer is swept over all layers and the best layer is
  reported; queries are filtered to items the model answers correctly under clean
  10-shot ICL. Produces both zero-shot and shuffled-10-shot curves.
\item
  \textbf{Fixed-layer full-label readout.} Teacher-forced accuracy of the full label
  at a fixed injection layer, on unfiltered queries.
\end{itemize}

Any cross-method comparison must hold the readout convention fixed; when quoting
numbers, state which convention they came from.

\subsection{Failure-mode hypotheses}\label{failure-mode-hypotheses}

For tasks with poor held-out zero-shot steering (synonym, word-length,
country--currency, partially antonym and product--company):

\begin{itemize}
\tightlist
\item
  \textbf{H1 — Overfitting to train tasks.} The isolation method (head set or basis)
  fits the train tasks and fails to transfer. \emph{Test:} task-specific fitting
  (Lever 3.2) should rescue the failing tasks if H1 holds.
\item
  \textbf{H2 — A single direction is not enough.} No single injected vector can
  induce these tasks zero-shot, regardless of how it is found. \emph{Test:} directly
  optimise an unconstrained \(v \in \mathbb{R}^{\dmodel}\) per task against the
  zero-shot objective; if even this upper bound fails, no head-selection method can
  succeed zero-shot on that task.
\item
  \textbf{H3 — Wrong success metric.} Zero-shot induction is simply the wrong bar
  for some tasks; the honest metric is uplift in a corrupted 10-shot context, where
  the format scaffold exists and steering supplies task identity. \emph{Test:} score
  the same isolation products under Lever 4.2/4.3; if the failures vanish there,
  adopt the mixed-context metric as primary.
\end{itemize}

These hypotheses are not exclusive: H2 and H3 can both be true (a task can need
context for the format \emph{and} more than one direction for the content).

\end{document}
